\documentclass[11pt]{article}

\usepackage[T1]{fontenc}
\usepackage{lmodern}
\usepackage{amsmath,amssymb,amsthm,mathtools}
\usepackage{geometry}
\usepackage{graphicx}
\usepackage{xcolor}
\usepackage{microtype}
\usepackage{enumitem}
\usepackage[hidelinks]{hyperref}

\geometry{margin=1in}
\setlength{\parindent}{0pt}
\setlength{\parskip}{0.65em}
\newcommand{\C}{\mathbb C}
\newcommand{\R}{\mathbb R}
\newcommand{\1}{\mathbf 1}
\newcommand{\norm}[1]{\left\lVert #1\right\rVert}
\newtheorem{theorem}{Theorem}

\title{Exact complex endpoint and \(C_0\) norms,\\
with strict interior phase loss}
\author{Candidate partial result for arXiv:2404.19540}
\date{9 August 2026}

\begin{document}
\maketitle

\begin{center}
\fcolorbox{black}{gray!10}{\parbox{0.91\textwidth}{
\textbf{Status: candidate partial result; verifier verdict: likely valid.}
The result closes both \(L^p\) endpoints and the adjacent \(C_0([0,1])\)
case for every complex order, and proves strict norm loss for every nonreal
order in the interior range.  It does not compute the square-space norm for
general \(1<p<\infty\).
}}
\end{center}

\section{Source question}

Alam, Chalendar, El Chami, Fricain, and Lef\`evre study the analytic
Riemann--Liouville semigroup in \cite{source}.  For \(\xi\in\C\) with
\(\operatorname{Re}\xi>0\), they write
\begin{equation}\label{eq:def}
 (V_\xi f)(x)=\frac{1}{\Gamma(\xi)}
 \int_0^x (x-u)^{\xi-1}f(u)\,du,
 \qquad 0\le x\le1.
\end{equation}
Their \(C_0([0,1])\) denotes the continuous functions on \([0,1]\) that
vanish at \(0\).

Remark~1 on PDF page~14 gives
\[
 \norm{V_\xi}_{\mathcal B(X)}
 \le \frac{1}{\operatorname{Re}(\xi)|\Gamma(\xi)|},
 \qquad X=C_0([0,1])\ \text{or}\ L^p(0,1),
\]
notes equality on \(L^1\) when \(\xi\) is real, and then states that the
exact \(L^p\)-norm is not known.  The full source passage is reproduced below.

\begin{figure}[ht]
  \centering
  \includegraphics[width=0.98\textwidth]{figures/open_problem_crop.png}
  \caption{Proposition~3.1(iv)--(vi) and Remark~1 of \cite{source}, PDF
  page~14.  The crop retains the complete exact-norm statement.}
  \label{fig:source}
\end{figure}

\clearpage
\section{Proof intuition}

The kernel modulus depends only on \(x-u\).  At the \(L^\infty\) endpoint,
one may align the phase of the single row at \(x=1\), so the row's \(L^1\)
mass is the exact norm.  At the \(L^1\) endpoint, normalized mass
concentrating at \(u=0\) makes the output converge to the first kernel column;
the column has the same \(L^1\) mass.  These are genuinely endpoint
mechanisms and do not produce an exact interior \(L^p\)-norm.

For the \(C_0\) restriction, the phase-aligning \(L^\infty\) function can be
cut off continuously near both endpoints.  The lost kernel mass tends to
zero.  More generally, the exact \(L^p\to C_0\) norm is simply the
\(L^{p'}\)-norm of the row at \(x=1\), because H\"older gives this as an upper
bound and evaluation at \(1\) supplies the matching lower bound.

In the interior range, compactness guarantees an extremizer.  If a nonreal
order attained the natural domination bound by its real-order counterpart,
equality in the triangle inequality would force the oscillatory phases
\((x-u)^{i\eta}\) to align simultaneously for almost every pair \((x,u)\).
Two distinct input points make that impossible as \(x\) varies.  This gives a
strict comparison, though not the still-unknown exact interior constant.

\section{Exact endpoint theorem}

Put \(\tau=\operatorname{Re}\xi\) and \(\eta=\operatorname{Im}\xi\).

\begin{theorem}\label{thm:main}
Let \(\xi=\tau+i\eta\in\C\) with \(\tau>0\), and let \(V_\xi\) be given by
\eqref{eq:def}.

\begin{enumerate}[label=\textup{(\roman*)}]
\item The exact square-space endpoint and \(C_0\) norms are
\begin{align*}
 \norm{V_\xi}_{L^1\to L^1}
 &=\norm{V_\xi}_{L^\infty\to L^\infty}
 =\norm{V_\xi}_{L^\infty\to C_0}
 =\norm{V_\xi}_{C_0\to C_0}\\
 &=\frac{1}{\tau|\Gamma(\xi)|}.
\end{align*}

\item Let \(1<p\le\infty\), let \(p'\) be conjugate to \(p\), and assume
\(\tau>1/p\).  Then
\begin{equation}\label{eq:lp-c0}
 \norm{V_\xi}_{L^p\to C_0}
 =\frac{1}{|\Gamma(\xi)|
 \bigl((\tau-1)p'+1\bigr)^{1/p'}}.
\end{equation}
For \(p=\infty\), the convention is \(p'=1\), so \eqref{eq:lp-c0}
reduces to \(1/(\tau|\Gamma(\xi)|)\).

\item The map \(V_\xi:L^1(0,1)\to C_0([0,1])\) is bounded if and only if
\(\tau\ge1\).  In that case
\begin{equation}\label{eq:l1-c0}
 \norm{V_\xi}_{L^1\to C_0}=\frac{1}{|\Gamma(\xi)|}.
\end{equation}
\end{enumerate}
\end{theorem}

\begin{proof}
For \(0<x\le1\), write the kernel row as
\[
 k_x(u)=\frac{\1_{(0,x)}(u)(x-u)^{\xi-1}}{\Gamma(\xi)}.
\]
Its modulus is
\begin{equation}\label{eq:modulus}
 |k_x(u)|=\frac{\1_{(0,x)}(u)(x-u)^{\tau-1}}
 {|\Gamma(\xi)|}.
\end{equation}

\emph{The \(L^p\to C_0\) row norm.}
Assume first that \(1<p\le\infty\) and \(\tau>1/p\).  H\"older's inequality
and \eqref{eq:modulus} give
\begin{align}
 |V_\xi f(x)|
 &\le \norm{k_x}_{p'}\norm{f}_p,\label{eq:holder}\\
 \norm{k_x}_{p'}
 &=\frac{x^{\tau-1/p}}
 {|\Gamma(\xi)|\bigl((\tau-1)p'+1\bigr)^{1/p'}}.
 \label{eq:row}
\end{align}
The exponent \(\tau-1/p\) is positive, so the right side is largest at
\(x=1\).  This proves the upper bound in \eqref{eq:lp-c0}.  Conversely,
evaluation at \(x=1\) shows that the operator norm is at least the norm of the
functional \(f\mapsto\int_0^1 k_1(u)f(u)\,du\), namely
\(\norm{k_1}_{p'}\).  Thus equality holds.

For completeness, the image really lies in the source's \(C_0\).  Extend
\(f\) and \(s\mapsto\1_{(0,1)}(s)s^{\xi-1}/\Gamma(\xi)\) by zero to \(\R\).
For finite \(p\), translation continuity of \(f\) in \(L^p(\R)\), followed
by H\"older, makes the convolution continuous.  For \(p=\infty\), use
translation continuity of the kernel in \(L^1(\R)\).  Its value at \(0\) is
zero by the Volterra support.

The same row argument at \(p=1\) uses the \(L^\infty\)-norm of \(k_x\).
It is uniformly finite exactly when \(\tau\ge1\), and then
\[
 \sup_{0<x\le1}\norm{k_x}_\infty=\frac1{|\Gamma(\xi)|}.
\]
This proves \eqref{eq:l1-c0}.  If \(\tau<1\), the functional associated to
\(k_1\) is unbounded on \(L^1\); inputs supported on intervals approaching
\(u=1\) give an explicit unbounded sequence.  When \(\tau\ge1\), continuity
of the output follows by translating the zero extension of \(f\) in
\(L^1(\R)\) against the bounded kernel.  This also covers the boundary
\(\tau=1\), even when \(\eta\ne0\).

\emph{The square \(L^1\)-norm.}
Tonelli's theorem yields, for every \(f\in L^1(0,1)\),
\begin{align*}
 \norm{V_\xi f}_1
 &\le\frac1{|\Gamma(\xi)|}
 \int_0^1 |f(u)|\int_u^1(x-u)^{\tau-1}\,dx\,du\\
 &=\frac1{\tau|\Gamma(\xi)|}
 \int_0^1(1-u)^\tau|f(u)|\,du\\
 &\le\frac1{\tau|\Gamma(\xi)|}\norm{f}_1.
\end{align*}
To obtain the reverse inequality, set
\[
 h(s)=\frac{\1_{(0,1)}(s)s^{\xi-1}}{\Gamma(\xi)}
 \quad(s\in\R),
 \qquad
 f_\varepsilon=\varepsilon^{-1}\1_{(0,\varepsilon)}.
\]
Then \(h\in L^1(\R)\), \(\norm{f_\varepsilon}_1=1\), and on \((0,1)\)
\[
 V_\xi f_\varepsilon
 =\frac1\varepsilon\int_0^\varepsilon h(\,\cdot-u\,)\,du.
\]
Translation continuity in \(L^1(\R)\) therefore gives
\(V_\xi f_\varepsilon\to h\) in \(L^1(0,1)\).  Consequently
\[
 \norm{V_\xi}_{L^1\to L^1}
 \ge\norm{h}_{L^1(0,1)}
 =\frac1{\tau|\Gamma(\xi)|},
\]
which proves the exact \(L^1\)-norm.

\emph{The \(L^\infty\) and \(C_0\) square norms.}
Taking \(p=\infty\) in the already proved row formula gives both
\[
 \norm{V_\xi}_{L^\infty\to C_0}
 =\norm{V_\xi}_{L^\infty\to L^\infty}
 =\frac1{\tau|\Gamma(\xi)|}.
\]
Restriction of the domain gives the same number as an upper bound on
\(C_0\to C_0\).  For the lower bound, choose continuous functions
\(0\le\chi_n\le1\), supported in \((0,1)\), with
\(\chi_n(u)\to1\) for every \(u\in(0,1)\), and define
\[
 f_n(u)=\chi_n(u)\frac{\Gamma(\xi)}{|\Gamma(\xi)|}
 (1-u)^{-i\eta}\quad(0\le u<1),
 \qquad f_n(1)=0.
\]
Because \(\chi_n\) vanishes near both endpoints, \(f_n\in C_0([0,1])\)
and \(\norm{f_n}_\infty\le1\).  The phases align at \(x=1\), so
\[
 V_\xi f_n(1)
 =\frac1{|\Gamma(\xi)|}
 \int_0^1(1-u)^{\tau-1}\chi_n(u)\,du
 \longrightarrow\frac1{\tau|\Gamma(\xi)|}
\]
by dominated convergence.  This is the matching lower bound on
\(C_0\to C_0\) and completes the proof.
\end{proof}

\section{Strict interior phase loss}

The full exact interior norm remains out of reach, but complex orders admit a
strict comparison with the corresponding positive real-order operator.

\begin{theorem}[Strict complex-to-real comparison]\label{thm:strict}
Let \(1<p<\infty\), \(\tau>0\), and \(\eta\in\R\setminus\{0\}\).  Then
\begin{equation}\label{eq:strict-comparison}
 \norm{V_{\tau+i\eta}}_{L^p\to L^p}
 <\frac{\Gamma(\tau)}{|\Gamma(\tau+i\eta)|}
 \norm{V_\tau}_{L^p\to L^p}.
\end{equation}
Moreover, for every \(\xi\in\C\) with \(\operatorname{Re}\xi=\tau>0\),
\begin{equation}\label{eq:strict-young}
 \norm{V_\xi}_{L^p\to L^p}
 <\frac1{\tau|\Gamma(\xi)|}.
\end{equation}
\end{theorem}

\begin{proof}
We first record compactness.  Extend
\[
 h_\xi(s)=\frac{\1_{(0,1)}(s)s^{\xi-1}}{\Gamma(\xi)}
\]
by zero to \(\R\).  Approximate \(h_\xi\) in \(L^1(\R)\) by continuous
kernels supported away from \(0\).  The corresponding integral operators on
\([0,1]\) have continuous kernels and are norm limits of finite-rank
operators, hence are compact.  Young's inequality transfers the
\(L^1\)-kernel approximation to operator norm.  Thus \(V_\xi\) is compact on
\(L^p(0,1)\).  Since \(L^p\) is reflexive, \(V_\xi\) attains its norm.

Pointwise domination gives
\begin{equation}\label{eq:pointwise-domination}
 |V_{\tau+i\eta}f(x)|
 \le \frac{\Gamma(\tau)}{|\Gamma(\tau+i\eta)|}
 V_\tau(|f|)(x).
\end{equation}
Suppose equality held in \eqref{eq:strict-comparison}, and let
\(f\in L^p\), \(\norm{f}_p=1\), attain the norm of
\(V_{\tau+i\eta}\).  Both inequalities in
\[
 \norm{V_{\tau+i\eta}f}_p
 \le\frac{\Gamma(\tau)}{|\Gamma(\tau+i\eta)|}
 \norm{V_\tau|f|}_p
 \le\frac{\Gamma(\tau)}{|\Gamma(\tau+i\eta)|}
 \norm{V_\tau}_{p\to p}
\]
would then be equalities.  In particular, equality holds in
\eqref{eq:pointwise-domination} for almost every \(x\).

Let \(E=\{u:f(u)\ne0\}\), which has positive measure, and write
\(\zeta(u)=f(u)/|f(u)|\) on \(E\).  Equality in the triangle inequality for
the defining integral implies that, for almost every \(x\), there is a
unimodular number \(\omega(x)\) such that
\begin{equation}\label{eq:phase-alignment}
 e^{i\eta\log(x-u)}\zeta(u)=\omega(x)
\end{equation}
for almost every \(u\in E\cap(0,x)\).  By Fubini, for almost every pair
\((u,v)\in E^2\) and almost every \(x>\max\{u,v\}\),
\[
 e^{i\eta[\log(x-u)-\log(x-v)]}=\frac{\zeta(v)}{\zeta(u)}.
\]
Choose such a pair with \(u\ne v\).  The left side is continuous in \(x\),
so equality almost everywhere on \((\max\{u,v\},1)\) makes it constant there.
Its logarithmic derivative is
\[
 i\eta\left(\frac1{x-u}-\frac1{x-v}\right),
\]
which is nonzero.  This contradiction proves
\eqref{eq:strict-comparison}.

It remains to prove \eqref{eq:strict-young} for real order.  The convolution
bound gives \(\norm{V_\tau}_{p\to p}\le1/\Gamma(\tau+1)\).  If equality
held, compactness would give a norming \(g\ge0\) with \(\norm g_p=1\).  Set
\[
 (S_sg)(x)=\1_{(s,1)}(x)g(x-s).
\]
Minkowski's inequality would give
\begin{align*}
 \norm{V_\tau g}_p
 &\le\frac1{\Gamma(\tau)}
 \int_0^1s^{\tau-1}\norm{S_sg}_p\,ds,\\
 \norm{S_sg}_p^p&=\int_0^{1-s}|g(u)|^p\,du.
\end{align*}
For \(s\) sufficiently close to \(1\), the latter norm is strictly less than
one.  Since the weight is positive, the integral is strictly less than
\(1/\Gamma(\tau+1)\), a contradiction.  Hence the real-order Young bound is
strict.  For nonreal order, combine this bound with
\eqref{eq:strict-comparison}; \eqref{eq:strict-young} follows in all cases.
\end{proof}

\section{What remains open}

Theorem~\ref{thm:main} settles the exact square norm for \(p=1,\infty\) and
the complete \(C_0\) analogue.  It also turns the source's
\(L^p\to C_0\) estimates into equalities.  The general exact value of
\(\norm{V_\xi}_{L^p\to L^p}\) for \(1<p<\infty\) remains open.

Theorem~\ref{thm:strict} shows that the imaginary part cannot be removed by a
phase choice: nonreal orders incur a strict loss relative to the real-order
domination constant.  It also proves that the elementary kernel bound is
never sharp in the interior.  But the positive real-order norm
\(\norm{V_\tau}_{L^p\to L^p}\) is itself unknown for general \(\tau\), so the
strict comparison does not yield an exact numerical value.

The endpoint constant cannot be extended to interior \(p\) merely by
interpolation: for the classical Volterra operator \(V_1\) on \(L^2(0,1)\),
the exact norm is \(2/\pi\), whereas
\(1/(\operatorname{Re}(1)|\Gamma(1)|)=1\).

\section{Verification and novelty status}

The proofs were checked line by line against the main likely failure points:
the complex phase in the \(L^1\) column limit, continuity into \(C_0\) at
\(\tau=1\), the endpoint cutoffs in the \(C_0\) lower bound, compactness and
norm attainment in the interior, and the Fubini phase-rigidity step.  All
checks pass.  The companion scripts numerically evaluate the explicit lower
bound families and a discretized strict-comparison example; they are only
sanity checks, not part of the proof.  Run them with
\begin{verbatim}
conda run --no-capture-output -n sandbox python \
  code/check_endpoint_norms.py
conda run --no-capture-output -n sandbox python \
  code/check_strict_phase_loss.py
\end{verbatim}

A bounded search covered the run indexes; the local parsed source series
arXiv:2108.02291, 2112.02669, 2408.10180, and 2410.00830; and arXiv/DOI web
queries for complex-order endpoint norms.  Carvalho-Neto and Fehlberg J\'unior
record the real-order \(L^1\) and \(L^\infty\) endpoint formula in
\cite{bochner}.  No complex-order endpoint or \(C_0\) formula was found.
The search also included strict complex-to-real norm comparison and
imaginary-order phase-loss phrases, with no exact match found.
Novelty confidence remains moderate because the proof is elementary and the
claim may be folklore.

\textbf{Human-review recommendation.}  Verify the endpoint norm-saturating
arguments and the Fubini phase-rigidity step, then perform a specialist
literature check before making any external originality claim.

\begin{thebibliography}{9}

\bibitem{source}
I.~Alam, I.~Chalendar, F.~El Chami, E.~Fricain, and P.~Lef\`evre,
\emph{Eventual Ideal Properties of the Riemann--Liouville Analytic Semigroup},
arXiv:2404.19540v2, 2026; Remark~1, p.~14.

\bibitem{bochner}
P.~M.~Carvalho-Neto and R.~Fehlberg J\'unior,
\emph{The Riemann--Liouville fractional integral in Bochner--Lebesgue spaces I},
arXiv:2108.02291v2, 2021; Theorem~4.1.

\end{thebibliography}

\end{document}
